%% ma: L12-B05-0003
%% lop: 12
%% chuong: 1
%% bai: 5
%% dang: 12.5.3
%% loai: TLN
%% muc: VD
%% dap_an: "49,6"
%% nguon: "(NBV)-ĐỀ SỐ 2-MỖI NGÀY 1 ĐỀ THI 2025.pdf (D02, MỖI NGÀY 1 ĐỀ - Đề số 2), Phần III Câu 2 (THPT Lương Tài số 2 - Bắc Ninh 2025)"
%% kien_thuc: "Bài 1 (cực trị của hàm số bậc ba), Bài 5 (ứng dụng đạo hàm vào bài toán thực tiễn); giải hệ phương trình bậc nhất hai ẩn"
%% ghi_chu: "Đề gốc ghi nhầm ngày 25/11/2014 (sửa thành 2024). Cách đánh số t: nếu 01/01 là t=0 thì 26/09 là t=269, hệ số không còn đẹp nhưng đáp số làm tròn vẫn 49,6 (kiểm bằng sympy: 49 616); lời giải theo t=270 như đáp án gốc. Đáp án gốc 49,6 đúng"
%% trang_thai: chinh-thuc
\begin{ex}
Trong khoảng thời gian từ ngày $01/01/2024$ đến hết ngày $30/12/2024$ nhóm nghiên cứu đã quan sát sự phát triển của một quần thể sinh vật $X$. Kết quả nghiên cứu chỉ ra rằng, tại ngày thứ $t$ của năm $2024$ (tính từ ngày $01/01/2024$) số cá thể sinh vật $X$ trong quần thể được ước lượng bởi hàm số $f(t)=-\dfrac1{300}t^3+bt^2+ct+12000$ (con), $0\le t\le365$ và ngày $26/09/2024$ là ngày có số lượng cá thể sinh vật $X$ nhiều nhất với $55\,740$ con. Ngày $25/11/2024$ số lượng cá thể sinh vật $X$ được ước lượng khoảng bao nhiêu nghìn con? (Kết quả làm tròn đến hàng phần chục).
\shortans{49,6}
\loigiai{Năm $2024$ là năm nhuận nên ngày $26/09$ là ngày thứ $270$ ($t=270$) và ngày $25/11$ là ngày thứ $330$ ($t=330$).
$f'(t)=-\dfrac1{100}t^2+2bt+c$. Hàm số đạt giá trị lớn nhất tại $t=270$ nên
$f'(270)=0\Leftrightarrow540b+c=729$ và $f(270)=55740\Leftrightarrow270b+c=405$.
Suy ra $b=\dfrac65$, $c=81$, $f(t)=-\dfrac1{300}t^3+\dfrac65t^2+81t+12000$. Thử lại: $f'(t)=-\dfrac{1}{100}(t-270)(t+30)$, $f'$ đổi dấu từ $+$ sang $-$ tại $t=270$.
Khi đó $f(330)=49\,620$ con $\approx49{,}6$ nghìn con.}
\end{ex}
